\specialchapter{术语索引}

绝对正则、正规（代数）：第~75 页。绝对正则、绝对正规代数：第~75 页。本质有限型代数：第~71 页。形式平展、形式非分歧代数：第~170 页, 习题 1 形式光滑代数：第~85 页。光滑代数：第~104 页。完全交环：第~65 页。Gorenstein 环：第~48 页。Grothendieck 环：第~169 页, 习题 5。局部整的环：第~16 页。Macaulay 环（或 Macaulay 环）：第~30 页。正规环：第~16 页。可表现环：第~58 页。正则环：第~55 页。Auslander-Buchsbaum（定理）：第~45 页。

Bass（定理）：第~49 页。Buchsbaum（模）：第~155 页, 习题 7。

Cohen（定理）：第~88 页。局部上同调（模）：第~142 页。完全割（理想）：第~62 页。对偶化复形：第~177 页, 习题 10 环的典范分支：第~15 页。

微差（理想）：第~167 页, 习题 7。环的同调维数：第~38 页。模的投射或内射维数：第~37 页。对偶化（复形）：第~177 页, 习题 10。对偶化（模）：第~125 页。Grothendieck 对偶：第~149 页。

本质有限型（代数）：第~71 页。本质生成（族）：第~71 页。

本质生成族：第~71 页。形式平展、形式非分歧（代数）：第~170 页, 习题 1。形式光滑（代数）：第~85 页。强割（子集）：第~28 页。

Gorenstein（环）：第~48 页。模的级：第~11 页。Grothendieck（环）：第~169 页, 习题 5。Grothendieck（对偶）：第~149 页。

Hartshorne（定理）：第~17 页。Hilbert-Burch（定理）：第~160 页, 习题 11。同调（维数）：第~38 页。完全割理想：第~62 页。内射（维数）：第~37 页。完全交（环）：第~65 页。光滑（形式光滑代数）：第~85 页。光滑（代数）：第~104 页。Macaulay（环）：第~30 页。Macaulay-Cohen（判别法）：第~31 页。Macaulay-Cohen（定理）：第~30 页。Macaulay（模）：第~30 页。Matlis（模）：第~110 页。Buchsbaum 模：第~155 页, 习题 7。局部上同调模：第~142 页。Matlis 模：第~110 页。对偶化模：第~125 页。Macaulay（或 Macaulay 的）模：第~23 页。完美模：第~158 页, 习题 4。纯模：第~27 页。正规（环）：第~16 页。完美（模）：第~158 页, 习题 4。强割子集：第~28 页。可表现（环）：第~58 页。模沿闭子集的深度：第~10 页。模、局部环的深度：第~2 页。投射（维数）：第~37 页。性质 \(S_k\)：第~151 页, 习题 8。素理想的 n 次符号幂：第~109 页。纯（模）：第~27 页。正则（环）：第~55 页。代数同态的提升：第~83 页。割（完全割理想）：第~62 页。Serre（定理）：第~20 页与第~54 页。模的底座：第~111 页。M-正则序列：第~8 页。Auslander-Buchsbaum 定理：第~45 页。Bass 定理：第~49 页。Cohen 定理：第~88 页。Hartshorne 定理：第~17 页。Hilbert-Burch 定理：第~160 页, 习题 11。Macaulay-Cohen 定理：第~30 页。Serre 定理：第~20 页与第~54 页。Zariski 定理：第~101 页与第~102 页。Zariski（定理）：第~101 页与第~102 页。
